% Generated draft: requires mathematical and source review.
\documentclass[11pt]{article}
\usepackage[margin=25mm]{geometry}
\usepackage{fontspec}
\setmainfont{Latin Modern Roman}
\usepackage{amsmath,amssymb,mathtools,unicode-math}
\setmathfont{Latin Modern Math}
\usepackage{enumitem,needspace,xurl,hyperref}
\setlength{\parindent}{0pt}
\setlength{\parskip}{5pt}
\setlength{\emergencystretch}{3em}
\newcommand{\N}{\mathbb{N}}
\newcommand{\Z}{\mathbb{Z}}
\newcommand{\Q}{\mathbb{Q}}
\newcommand{\R}{\mathbb{R}}
\newcommand{\C}{\mathbb{C}}
\newcommand{\F}{\mathbb{F}}
\newcommand{\A}{\mathbb{A}}
\newcommand{\PP}{\mathbb{P}}
\newcommand{\E}{\mathbb{E}}
\newcommand{\eps}{\varepsilon}
\newcommand{\dd}{\,\mathrm{d}}
\begin{document}
\section*{UnsolvedMath ID: 30005786}
\textbf{Source problem code:} catalogue entry
\subsection{Definitions and mathematical statement}
The supplied entry, titled ``Definition of the Coxeter Galaxy,'' contains the question ``What is the Coxeter galaxy?'' It supplies neither a definition of the Coxeter galaxy nor a mathematical assertion about it. In particular, the underlying objects, incidence relations, scope, and any intended geometric structure are unspecified.

The supplied statement-verification note reports that this question is a section heading in the original report, followed there by a definition, and is not an open problem. That definition is absent from the supplied record. Consequently, no precise definition or quantified open-problem statement can be recovered from the supplied material. A definition or conjecture about the Coxeter galaxy must not be inferred from this heading alone.

The registry labels the entry \texttt{partially\_solved}, while its statement status is \texttt{unrecoverable}. These are supplied curation labels. The background mentions global geometry and an associated isomorphism problem, but gives no precise formulation of either; neither can serve as a reconstructed mathematical target.
\subsection{Short English statement}
This entry appears to preserve a section heading asking for the definition of the Coxeter galaxy. The supplied verification note says it is not an open problem, and the definition itself is missing. A precise mathematical statement therefore requires additional source text and human review.
\subsection{Sources}
\begin{itemize}
\item Local UnsolvedMath registry: supplied record with permanent ID \texttt{30005786}, exact source code \texttt{catalogue entry}, in the collection ``Oberwolfach Reports Open Problems.'' The supplied proposed year is 2024; no proposer is named.
\item Supplied report citation: \emph{Mini-Worskhop: Artin Groups meet Triangulated Categories} (2024), Oberwolfach Reports, pp.\ 203--234, DOI \texttt{10.4171/owr/2024/4}. Source URL and literature-source link labeled ``original OWR report'': \url{https://doi.org/10.4171/owr/2024/4}.
\item Supplied literature source: Santos Rego--Schwer, \emph{The galaxy of Coxeter groups}, \url{https://doi.org/10.1016/j.jalgebra.2023.12.006}.
\item Literature assessment dated \texttt{2026-08-22}: the supplied curation labels the topic partially solved and states that a rigorous definition, local finiteness, infinite dimensionality, finite connected components, and low-rank structure are known, while full global geometry and an associated isomorphism problem remain open. This is supplied dated curation derived from the OWR material, not an independently verified assessment. Its time-sensitive status and mathematical claims require independent verification.
\item The supplied statement says that the report is based on reference [1], dedicated to the memory of Jacques Tits for his influential work on Coxeter groups. It attributes the terms ``Coxeter group'' and ``Coxeter diagram'' to Tits acting as an honorary Bourbaki. These are supplied historical attributions, not independently verified here. The record does not explicitly identify reference [1].
\end{itemize}
\end{document}
